\section{Introduction}\label{sec:intro}

A controlling shareholder who wants to sell a division of a listed company has, in Korea, a way of doing it in two steps. The company first places the division in a wholly owned subsidiary, a transaction the Korea Exchange counted \KLKrxCaptive{} times among \KLKrxTotal{} divisions between 1998 and 2022 and that \citeA{kang2025} call a captive spin-off; the parent's shareholders vote on this step, and nothing in the consolidated accounts changes. Then, months or years later, the parent's board sells part or all of the subsidiary or lists it, and from this step the parent's shareholders are excluded \cite{kang2025}. The two steps are what the public controversy of 2020--2022 was about. LG Chem's battery business was carved out in September 2020 and listed in January 2022, and over that interval LG Chem's own stock lost more than \KLLGDrop{} percent of its value \cite{kang2025}. By February 2022 the English-language press was asking why Korean-style divisions ``wreak havoc on minor shareholders'' \cite{kang2025}. In September 2022 the Financial Services Commission answered with three measures: a mandatory disclosure of the division's purpose and of the measures taken to protect shareholders, an appraisal right for dissenting shareholders, and an exchange review of any listing of the subsidiary within five years \cite{fsc2022a}.

The economics of the two steps is the economics of the principal--principal conflict \cite{young2008}. The controller's gain from the second step---proceeds that fund the group, a stake placed with an investor on terms the controller sets, a listing that gives the subsidiary its own currency---does not depend on how much of the parent the controller owns, while the controller's share of the parent's loss does. The literature on Korean business groups has documented this arithmetic in mergers \cite{bae2002}, in private securities offerings \cite{baek2006} and in intragroup support \cite{bae2008}. \citeA{kang2025} have now documented it for captive spin-offs. Across \KLN{} divisions announced between 2001 and 2021, the parent's stock earns a small positive abnormal return at the announcement and an abnormal return of \KLBhar{} percent against a matched portfolio over the following year; the loss is concentrated where the controlling owner's stake is small, and in the divisions whose subsidiary was eventually disposed of. They close with a sentence they do not test: that investors are, in most cases, unlikely to discern the controlling owner's intention at the announcement.

We test it. The division filing is a long structured document. The parent must state in it the purpose of the division, the business and capital of the new company, whether the new company will apply for listing and, from October 2022, the measures it takes to protect shareholders. A controller who intends to sell or list the subsidiary has ordinary commercial reasons to say things in it that a controller who intends to keep the subsidiary does not. We have two large language models read the division filings of \NScored{} captive spin-offs announced between 2001 and 2026, each as of its date, and forecast whether the parent will still own all of the new company three to five years later. We then trace what the parent actually did through its subsequent annual reports. The reading predicts the outcome. Across the \NGateOne{} divisions whose outcome is known, the two-model forecast orders disposed from retained subsidiaries with an area under the ROC curve of \AUCPooledLLM{}, against \AUCPooledFin{} for a model on the financial characteristics stated in the same filing. It does nearly as well when the filings are stripped of every name and date, so it is reading rather than remembering. It is sharpest, at \AUCEarlyPooled{}, about the disposals that come within a year of the division.

The market does not read the filing this way. Across the listed parents the three-day announcement return rises with the disposal forecast, by \AnnCarSlopePP{} percentage points from a filing that reads as a certain retention to one that reads as a certain disposal. The market-adjusted return over the following year falls with it: a Spearman correlation of \SpearmanScoreBhar{} across \NScoreOneY{} parents, a within-year rank coefficient of \RankSpecThreeCoef{} after size, book-to-market, momentum, market and filing controls, and \QQWithinYear{} percentage points between the top and bottom quintiles of the forecast. The loss is realised where the forecast said it would be. The one-year return of parents whose subsidiary changed hands within a fiscal year of the division differed from that of parents that kept it by \EarlyNeverDiff{} percentage points; parents whose disposal came later did not trail within the year; and once realised early disposal is in the regression, the forecast itself adds nothing. The chain runs from the controller's stake. Parents whose controlling block before the announcement was in the lowest tercile (up to \LowBlockMax{} percent of common shares) filed divisions that read more strongly as disposals, disposed more often (\LowBlockDisp{} against \HighBlockDisp{} percent in the top tercile) and earned a median one-year return of \LowBlockMedian{} percent against \HighBlockMedian{} percent. The forecast falls with the block (\ScoreOnBlock{} per unit of the block). After the 2022 reform the filings changed within weeks: shareholder-protection statements rose from \ProtectPre{} to \ProtectPost{} percent of filings and the announcement reaction from \CarPre{} to \CarPost{} percent. The disposal forecast and the rate of early disposal did not move over the horizon we can observe.

These results bear on two conversations in this journal. The first is the review of business-group governance by \citeA{aguilera2025} and the studies of controlling-owner expropriation and its restraints that surround it \cite{shi2022, bian2023, huang2025, na2025, qian2026}. The open question there is less whether controllers extract than whether anything short of ownership restrains them. A captive spin-off is a clean case for that question because the extraction is deferred and discretionary: the controller announces a reorganisation, obtains a vote on it, and keeps the option to realise the transfer later. We show that the option was legible in the controller's own disclosure at the time of the vote; that minority investors priced the disclosure in the wrong direction; and that a reform which made the disclosure mandatory changed the disclosure and its reception before it changed anything the disclosure was about. The second conversation is about disclosure as a governance instrument \cite{he2025, karamatzanis2026}. Our reading of the filings is a measure of what disclosure contained before anyone was required to put it there, and the 2022 episode shows what mandating it added.

Does the controlling owner's intention to dispose of a carved-out subsidiary show in the division filing before minority shareholders bear the loss, and does the market price it? Section~\ref{sec:setting} sets out the captive spin-off, the 2022 reform and the four hypotheses. Section~\ref{sec:data} describes the filings, the outcome tracing, the disposal signal and the return and ownership measures. Section~\ref{sec:results} reports the results, Section~\ref{sec:discussion} maps them back to the theory and to the tests of \citeA{kang2025} and draws out the implications, and Section~\ref{sec:conclusion} states what the evidence cannot support.
